\begin{enumerate}[label=\textbf{1.1.\arabic*.},ref=1.1.\arabic*]
\setlength{\itemsep}{0.35em}
\item \label{formula-1-1-1} \textbf{Tangents and radii.} The radius to a tangent point is perpendicular to the tangent. Tangents from one external point have equal lengths. Thus $OP^2=OT^2+PT^2$ for circle radius $OT$.
\item \label{formula-1-1-2} \textbf{Coordinate distance and circle equation.}
\begin{equation*}\tag{1.1.2.1}
AB=\sqrt{(x_2-x_1)^2+(y_2-y_1)^2},\qquad
(x-h)^2+(y-k)^2=r^2
\end{equation*}
for a circle with centre $(h,k)$ and radius $r$.
\item \label{formula-1-1-3} \textbf{Pairs of linear equations.} For $a_1x+b_1y+c_1=0$ and $a_2x+b_2y+c_2=0$, the pair has a unique solution when $\frac{a_1}{a_2}\ne\frac{b_1}{b_2}$, infinitely many solutions when all three ratios are equal, and is inconsistent when
\begin{equation*}\tag{1.1.3.1}
\frac{a_1}{a_2}=\frac{b_1}{b_2}\ne\frac{c_1}{c_2}.
\end{equation*}
\item \label{formula-1-1-4} \textbf{Probability.} $P(E)=\frac{n(E)}{n(S)}$ for equally likely outcomes; for repeated independent trials, $P(\text{at least one }E)=1-(1-P(E))^n$.
\item \label{formula-1-1-5} \textbf{Quadratic equations.} The discriminant is $\Delta=b^2-4ac$; equal roots occur when $\Delta=0$, and the roots have sum $-b/a$ and product $c/a$.
\item \label{formula-1-1-6} \textbf{Similar triangles.} Corresponding sides of similar triangles are proportional. If a line parallel to one side of a triangle meets the other two sides, it divides them proportionally: $\frac{AD}{DB}=\frac{AE}{EC}$.
\item \label{formula-1-1-7} \textbf{Arc length.} For central angle $\theta$ in degrees, $s=\frac{\theta}{360^\circ}2\pi r$. In radians, $s=r\theta$.
\item \label{formula-1-1-8} \textbf{Arithmetic progressions.} For first term $a$ and common difference $d$,
\begin{equation*}\tag{1.1.8.1}
a_n=a+(n-1)d,\qquad S_n=\frac n2\bigl(2a+(n-1)d\bigr)=\frac n2(a+a_n).
\end{equation*}
\item \label{formula-1-1-9} \textbf{HCF and LCM.} $\operatorname{HCF}(a,b)\operatorname{LCM}(a,b)=ab$; common divisors are found using prime factorization or Euclid's algorithm.
\item \label{formula-1-1-10} \textbf{Section formula.} If $AP:PB=m:n$, then $P=\left(\frac{nx_1+mx_2}{m+n},\frac{ny_1+my_2}{m+n}\right)$ for $A(x_1,y_1)$, $B(x_2,y_2)$.
\item \label{formula-1-1-11} \textbf{Heights and distances.} In a right triangle, $\tan\theta=\frac{\text{vertical height}}{\text{horizontal distance}}$ and $\sin\theta=\frac{\text{vertical height}}{\text{string length}}$.
\end{enumerate}
